Lo spazio per la chiave e' strettamente legato al modello di ME
che si sta analizzando.
Le differenze, in base alla modifica, possono essere anche ordini di
grandezza piu' grandi.

Alcune assunzioni fatte sono poi veritiere solo dal 1926-1931 quando ancora
non si conoscieva "quasi nulla" della macchina cifrante.


\subsubsection{Cavi scambianti}
Le combinazioni sono analoghe a "quante C coppie distinte posso formare da un alfabeto di 26 lettere?".
Siccome ci sono 26 lettere totali, il numeri di cavi scambiati sara' compreso:
$0 \le cavi \le 13$ dove 13 e' semplicemente 26 lettere diviso due.

Calcoliamo il numero di possibili combinazioni dal seguente binomio di newton:

$$\binom{26}{2c} = \frac{26!}{(26-2c)! \times c! \times 2^{c}}$$
Dove con "c" intendiamo il numero di cavi.

Si nota con un conto esausitovo che:
$$max(\binom{26}{2c}} \longrightarrow  c = 11$$

Infatti:
$$\binom{26}{22} = 205552193096250$$

\subsubsection{Cavi dentro i rotori}
Le combinazioni in questo caso sono analoghe a "quanti modi ho di diporre 26 lettere?".
In questo caso il numero sarebbe: $26!$.

Questo numero di fatto sara' solo teorico.
I dipartimenti avevano in dotazione dei dischi standard.
L'esercito non fini mai ad usare tutti, o parte, dei $26!$ dischi possibili.

\subsubsection{Posizione iniziale anelli}
Siccome parliamo di elementi distinti possiamo calcolarne uno e da li' capire per il numero di rotori.
Un anello poteva prendere 26 posizioni possibili.
Quindi, con 3 rotori potevano contribuire con $26^{3} = 17576$ posizioni totali.
Piu' avanti con 4 rotori si raggiunse: $26^{4}= 456976$ combinazioni totali.

Alcune macchine potevano avere anche due tipi di anelli, ma ai fini del conto delle possiblita'
questo concetto non ha effetti pratici.

\subsubsection{Possibili riflessori}
Analogamente ai cavi scambianti:
"Quante possibili 13 coppie distinte posso generare da 26 lettere?".

$$ \frac{26!}{13!\times2^13}= 7905853580625$$

Come specificato in "cavi dentro i rotori", di fatto ogni reparto aveva il suo riflessore.

\subsubsection{Scelta rotori}
Durante la guerra la ME passo' dall'avere i soli 3 rotori di base fino ad un
massimo di 8 rotori.
Solo tre rotori tuttavia potevano essere alloggiati dentro la ME.
Parte della chiave quindi divenne anche scegliere quale rotore finiva dentro la ME.

$$ \binom{n}{3}$$
dove: $ 3 \le n \le 8$ in base al tipo di modello.
